\documentclass[10pt,a4paper]{amsart}
\usepackage{amsmath,amssymb,amsthm,graphicx}
% Retained minimal preamble for the Stacks Project chapter more-algebra.tex (More on Algebra).
% The chapter file pinned-original-source.tex carries only \input{preamble} on its first line, so no
% preamble text lives inside the chapter itself.  The Stacks Project uses one single shared preamble for
% every chapter; the snapshot stages it identically as sources/stacks-chow/preamble.tex and
% sources/stacks-etale-cohomology/preamble.tex.  The block below is copied mechanically, line for line,
% from sources/stacks-chow/preamble.tex lines 14-18 (the Xy-pic setup) and 154-217 (the native theorem
% declarations together with the mathematical macros).  Nothing was added, removed or redefined; the
% wrapper class and the amsmath/amssymb/amsthm/graphicx base packages are supplied by the assembler.
% Chapter: more-algebra.tex -- More on Algebra.
% For commutative diagrams we use Xy-pic
\usepackage[all]{xy}

% We use 2cell for 2-commutative diagrams.
\xyoption{2cell}

% Theorem environments.
%
\theoremstyle{plain}
\newtheorem{theorem}[subsection]{Theorem}
\newtheorem{proposition}[subsection]{Proposition}
\newtheorem{lemma}[subsection]{Lemma}

\theoremstyle{definition}
\newtheorem{definition}[subsection]{Definition}
\newtheorem{example}[subsection]{Example}
\newtheorem{exercise}[subsection]{Exercise}
\newtheorem{situation}[subsection]{Situation}

\theoremstyle{remark}
\newtheorem{remark}[subsection]{Remark}
\newtheorem{remarks}[subsection]{Remarks}

\numberwithin{equation}{subsection}

% Macros
%
\def\lim{\mathop{\mathrm{lim}}\nolimits}
\def\colim{\mathop{\mathrm{colim}}\nolimits}
\def\Spec{\mathop{\mathrm{Spec}}}
\def\Hom{\mathop{\mathrm{Hom}}\nolimits}
\def\Ext{\mathop{\mathrm{Ext}}\nolimits}
\def\SheafHom{\mathop{\mathcal{H}\!\mathit{om}}\nolimits}
\def\SheafExt{\mathop{\mathcal{E}\!\mathit{xt}}\nolimits}
\def\Sch{\mathit{Sch}}
\def\Mor{\mathop{\mathrm{Mor}}\nolimits}
\def\Ob{\mathop{\mathrm{Ob}}\nolimits}
\def\Sh{\mathop{\mathit{Sh}}\nolimits}
\def\NL{\mathop{N\!L}\nolimits}
\def\CH{\mathop{\mathrm{CH}}\nolimits}
\def\proetale{{pro\text{-}\acute{e}tale}}
\def\etale{{\acute{e}tale}}
\def\QCoh{\mathit{QCoh}}
\def\Ker{\mathop{\mathrm{Ker}}}
\def\Im{\mathop{\mathrm{Im}}}
\def\Coker{\mathop{\mathrm{Coker}}}
\def\Coim{\mathop{\mathrm{Coim}}}

% Boxtimes
%
\DeclareMathSymbol{\boxtimes}{\mathbin}{AMSa}{"02}

%
% Macros for moduli stacks/spaces
%
\def\QCohstack{\mathcal{QC}\!\mathit{oh}}
\def\Cohstack{\mathcal{C}\!\mathit{oh}}
\def\Spacesstack{\mathcal{S}\!\mathit{paces}}
\def\Quotfunctor{\mathrm{Quot}}
\def\Hilbfunctor{\mathrm{Hilb}}
\def\Curvesstack{\mathcal{C}\!\mathit{urves}}
\def\Polarizedstack{\mathcal{P}\!\mathit{olarized}}
\def\Complexesstack{\mathcal{C}\!\mathit{omplexes}}
% \Pic is the operator that assigns to X its picard group, usage \Pic(X)
% \Picardstack_{X/B} denotes the Picard stack of X over B
% \Picardfunctor_{X/B} denotes the Picard functor of X over B
\def\Pic{\mathop{\mathrm{Pic}}\nolimits}
\def\Picardstack{\mathcal{P}\!\mathit{ic}}
\def\Picardfunctor{\mathrm{Pic}}
\def\Deformationcategory{\mathcal{D}\!\mathit{ef}}
\usepackage{cleveref}
\crefname{originallemma}{original lemma}{original lemmas}
\Crefname{originallemma}{Original lemma}{Original lemmas}
\makeatletter
\newlabel{more-algebra-lemma-baire-category-complete-module}{{0CQU}{source}{}{sourcepointer-1}{}}
\newlabel{more-algebra-lemma-baire-category-complete-module@cref}{{[originallemma][0][]0CQU}{[1][0][]source}{}{}{}}
\makeatother
\title{Baire category theorem for a complete linearly topologized abelian group with a countable fundamental system}
\author{The Stacks Project Authors}
\date{stacks-project:more-algebra.tex:local-snapshot-20261012; GNU FDL 1.2 or later}
\begin{document}
\maketitle
\setcounter{subsection}{1}
\noindent\textit{Editorial note.} All delivered mathematical text is the author's own, copied from the pinned Stacks Project chapter file more-algebra.tex; no statement or proof step was rewritten, reordered or removed.\par\medskip
\noindent\textit{Editorial note.} The chapter file pinned-original-source.tex carries only \textbackslash{}input\{preamble\} on its first line, so no preamble text lives inside the chapter. The retained file sources/stacks-more-algebra/retained-preamble.tex is a mechanical, line-for-line copy of the project's single shared preamble as staged in the snapshot, namely sources/stacks-chow/preamble.tex lines 14-18 (Xy-pic setup) and 154-217 (theorem declarations and mathematical macros). Nothing was added, removed or redefined, no package was added, and the excerpt serves only as preamble, never as a span coordinate.\par\medskip
\noindent\textit{Editorial note.} Both counted spans start at the \textbackslash{}begin line of their environment, so the statement and the author's proof are delivered whole; surrounding prose of the section is deliberately outside the counted spans and is carried as standard prerequisites instead.\par\medskip
\noindent\textit{Editorial note.} The citation fields (source\_id, source\_version, authors, source\_url, license.name, source\_title) are identical to those of the cards already published from this chapter, so that the batch selector treats them as one source instead of deferring them as a same-source clash.\par\medskip
\noindent\textit{Allowed standard prerequisites (editorial).}\begin{itemize}
\item Linearly topologized abelian group with a countable fundamental system of neighbourhoods of zero. definition discussion of the same section, source lines 8996-9022, before the counted statement. Supplies the nested neighbourhood basis M\_n and the meaning of completeness used to take the limit.
\item Completeness of a topological abelian group in terms of Cauchy sequences. used at proof lines 9067-9069 to produce the limit y of the constructed sequence. Turns the constructed sequence, whose successive differences lie in the shrinking neighbourhoods, into an actual element y of M.
\item A subset is dense if and only if it meets every nonempty open subset. used at proof lines 9057-9058 and 9061-9062 when the y\_i are chosen inside the open cosets x + M\_k and y\_1 + M\_\{k\_1\}. Provides the points y\_i in U\_i at every step of the induction.
\end{itemize}\medskip
\ifdefined\hypertarget\hypertarget{sourcepointer-1}{}\fi
\noindent\textit{Source reference.} Original lemma 0CQU: focused-seeds/AGP/sources/stacks-snapshot-20261011/tags.txt, source lines 15974--15974. The lemma's own canonical tag, quoted so that readers can resolve the label inside the delivered text; it points at the counted statement itself and carries no proof content.\par\medskip
\begin{lemma}[Baire category theorem]
\label{lemma-baire-category-complete-module}
Let $M$ be a topological abelian group. Assume $M$ is linearly
topologized, complete, and has a countable fundamental system of
neighbourhoods of $0$. If $U_n \subset M$, $n \geq 1$
are open dense subsets, then $\bigcap_{n \geq 1} U_n$ is dense.
\end{lemma}

\begin{proof}
Let $U_n$ be as in the statement of the lemma. After replacing $U_n$ by
$U_1 \cap \ldots \cap U_n$, we may assume that
$U_1 \supset U_2 \supset \ldots$.
Let $M_n$, $n \in \mathbf{N}$ be a fundamental
system of neighbourhoods of $0$. We may assume that
$M_{n + 1} \subset M_n$. Pick $x \in M$. We will show that
for every $k \geq 1$ there exists a $y \in \bigcap_{n \geq 1} U_n$
with $x - y \in M_k$.

\medskip\noindent
To construct $y$ we argue as follows.
First, we pick a $y_1 \in U_1$ with $y_1 \in x + M_k$.
This is possible because $U_1$ is dense and $x + M_k$ is open.
Then we pick a $k_1 > k$ such that $y_1 + M_{k_1} \subset U_1$.
This is possible because $U_1$ is open.
Next, we pick a $y_2 \in U_2$ with $y_2 \in y_1 + M_{k_1}$.
This is possible because $U_2$ is dense and $y_2 + M_{k_1}$ is open.
Then we pick a $k_2 > k_1$ such that $y_2 + M_{k_2} \subset U_2$.
This is possible because $U_2$ is open.

\medskip\noindent
Continuing in this fashion we get a converging sequence
$y_i$ of elements of $M$ with limit $y$. By construction
$x - y \in M_k$. Since
$$
y - y_i = (y_{i + 1} - y_i) + (y_{i + 2} - y_{i + 1}) + \ldots
$$
is in $M_{k_i}$ we see that $y \in y_i + M_{k_i} \subset U_i$
for all $i$ as desired.
\end{proof}

\end{document}
